\documentclass[11pt]{article}
\usepackage[T1]{fontenc}
\usepackage{lmodern,amsmath,amssymb,microtype}
\usepackage[margin=1.02in]{geometry}
\usepackage[colorlinks=true,urlcolor=blue,citecolor=blue]{hyperref}
\setlength{\emergencystretch}{2em}
\hypersetup{pdfauthor={Wentao Li},pdftitle={Frucht's theorem and weak forms of anti-foundation}}
\title{Frucht's theorem and weak forms of anti-foundation}
\author{Wentao Li\\[2pt]\small Prepared with assistance from OpenAI Codex}
\date{November 2026}
\begin{document}
\maketitle
\begin{abstract}
Over ZF with Foundation deleted, Frucht's theorem implies that every extensional accessible pointed digraph has an extensional, isomorphism-extensional representative. This reverses the entire chain of three implications asked about by Chen and Ju. The proof supplies the graph labels for their construction from a single set-sized application of Frucht's theorem.
\end{abstract}

\paragraph{Statement.} Write FT for the assertion that every group is the automorphism group of a simple graph. A pointed digraph is accessible if every vertex is reached by a finite directed path from its point. A decoration is a root-preserving map taking each successor set onto the successor set of the image. Extensionality means distinct vertices have different successor sets. Isomorphism-extensionality means their rooted descendant graphs are nonisomorphic. The principle $\mathrm{CWFAFA}_2$ asserts that every extensional accessible pointed digraph receives a decoration from an accessible pointed digraph having both properties. Restricting targets to canonical membership pictures gives $\mathrm{WFAFA}_2$. Dropping extensionality of the source in the latter gives $\mathrm{WFAFA}'_2$.

\paragraph{Theorem.} In $\mathrm{ZF}-\mathrm{Foundation}$,
\[
 \mathrm{FT}\ \longleftrightarrow\ \mathrm{CWFAFA}_2\ \longleftrightarrow\
 \mathrm{WFAFA}_2\ \longleftrightarrow\ \mathrm{WFAFA}'_2 .
\]
Chen and Ju prove the implications from left to right starting at $\mathrm{CWFAFA}_2$, and $\mathrm{WFAFA}'_2\Rightarrow\mathrm{FT}$. We prove $\mathrm{FT}\Rightarrow\mathrm{CWFAFA}_2$.

\paragraph{A set of reusable labels.} For any set $V$, FT furnishes an indexed family $(R_v)_{v\in V}$ of pairwise nonisomorphic rigid connected pointed simple graphs (Chen and Ju, Theorem~1). Form their tagged disjoint union, introduce a new point $r$, and join $r$ to just the specified point of each piece. The resulting pointed graph $S$ is rigid. Indeed, removing $r$ recovers the components and their distinguished attachment points. Their types force each component to be preserved and rigidity then fixes it pointwise. The attachment points give an injection $i:V\to S$.

Form a set $U$ of finite expressions starting with distinct leaf symbols for the elements of $S$ and for natural numbers, and a constructor $[F]$ for each finite set $F$ of already constructed expressions. Take the union over finite construction depths. Technically leaf and constructor records use distinct tags. Equality of constructor records is equality of finite sets of children. This is an ordinary recursion on $\omega$ even if the external sets indexing the leaves are ill-founded. The constructor codes ordered pairs by the Kuratowski formula and hence codes tagged tuples and unordered pairs within $U$. We use these codes below. Regard $U$ as a structure with the two leaf sorts, $S$-adjacency and its point on the first sort, zero and successor on the second, and the child relation from constructors to their arguments. Every automorphism is trivial, by rigidity on the leaves and induction on expression depth.

For every $x\in U$ there is a uniformly specified unpointed connected simple graph $C_x$, with at least two vertices, all its vertices coded in $U$, such that $C_x\cong C_y$ implies $x=y$. Here are details ensuring that the vertex codes require no further choices. Add a unary predicate marking $x$ to the preceding fixed finite signature. Make one coloured node per element, and for each true relation tuple make a relation node and an occurrence node for each position, joining the occurrence to the relation node and the appropriate element. Use relation and position colours. Add a unique anchor adjacent to all these nodes. Subdivide every edge once with a new colour. The resulting connected skeleton is bipartite. At each skeleton vertex attach a pendant odd cycle whose length encodes its colour. Each odd simple cycle is a marker, and its attachment is its unique vertex with a neighbour outside that cycle. Thus isomorphisms recover the coloured skeleton, then the marked structure. Every node, including each subdivision and cycle position, has a distinct finite tagged record over $U$ and fixed integers. In particular the assignment $x\mapsto C_x$ is a function produced by Separation and Replacement. We identify the target vertices with their codes via $i$.

\paragraph{The expansion.} Let $G$ be the target. Let $W$ be its vertices with well-founded descendant relation, using the subset definition of well-foundedness, and put $I=V(G)\setminus W$. There is no edge from $W$ to $I$. Every $u\in I$ has a successor in $I$: if all successors are well-founded, so is their union below $u$ (equivalently apply well-founded ranks to the successors). If $I$ is empty, well-founded extensional collapse already gives the required property. Otherwise the root lies in $I$ and every vertex of $I$ is accessible there.

Use the $U$ just built for $V(G)$. Start with one tagged root $x$ mapped to the root of $G$. Expand each newly introduced vertex once, in successive finite stages. For $u=\pi(x)$ and every successor $c$ of $u$ in $I$, introduce private children $b(x,c)$ and $d(x,c,q)$ for all $q\in V(C_x)$, mapping them all to $c$. For each edge $e$ of $C_x$ and each successor $z$ of $c$ in $I$, add $s(x,c,e,z)$ mapping to $z$, with incoming edges from precisely the two $d(x,c,q)$ with $q\in e$. These are distinct tagged records, hence again elements of the original $U$. The root uses its own tag. Include the $s$-vertices among the new vertices to be expanded. Finite-stage recursion and union provide a set $K$. All vertices are accessible, all edges respect the projection, and the $b$-vertices ensure the successor-surjectivity required of a decoration. Each vertex has a child of indegree one identifying it as parent, so $K$ is extensional.

At a vertex $x$, join different children in an auxiliary simple graph if they share a child, and discard isolated vertices. The only vertices with two parents are the $s$-vertices. An adjacent pair in the auxiliary graph consists of two $d$-vertices made for this very $x$ and the same $c$, since each $d$ has just its named parent. Pre-existing children of $x$ when $x$ is itself a $d$ are $s$-vertices and hence isolated in this test. For each $c$ the auxiliary edges are exactly those of $C_x$, since $c$ has a successor and $C_x$ is connected of order at least two. Thus the auxiliary components are a nonempty family of copies of $C_x$. This construction is intrinsic to the rooted descendant graph. An isomorphism between the descendants rooted at $x$ and $y$ gives $C_x\cong C_y$, whence $x=y$. This decoding and the expansion are those of Chen and Ju's Theorem~17. The reusable set of labels above removes its global or dependent-choice hypothesis.

\paragraph{Restoring well-founded vertices.} Adjoin a single tagged copy of $W$ with its original edges. Add $x\to w$ exactly when $\pi(x)\to w$ in $G$. Extending $\pi$ by the identity on $W$ gives a decoration. Finite paths in $I$ lift using the private $b$-children, so every restored vertex is accessible. Each $K$-vertex still has a private child and its $K$-descendants form a nonempty subset without a terminal member. The well-founded vertices are therefore intrinsically exactly $W$. Isomorphic rooted descendant graphs at $K$-vertices restrict on deleting $W$ to isomorphic rooted graphs in $K$. Within $W$ the usual well-founded extensional collapse distinguishes vertices. This proves both required extensionality properties. All recursions after the initial choice of one indexed family were defined by functions.\hfill$\square$

\paragraph{Source and priority.} J. Chen and D. Ju, \emph{Frucht's theorem and other set-theoretic principles below the axiom of choice and the axiom of foundation}, §§2, 7--8, \url{https://arxiv.org/abs/2607.23891v2}. The equivalence above addresses their second open-question bullet. The stronger injective representation principle $\mathrm{FAFA}_2$ is outside the claim. This is a research draft pending independent proof and priority review.
\end{document}
